\documentclass[10pt,a4paper]{amsart}
\usepackage[margin=19mm]{geometry}
\usepackage{amsmath,amssymb,mathtools}
\usepackage[scr=rsfs,cal=euler]{mathalfa}
\usepackage{xfrac}
\usepackage[unicode,hidelinks,bookmarks=false]{hyperref}
\newcommand{\ie}{i.e.\ }
\newcommand{\cf}{cf.\ }
\newcommand{\resp}{respectively\ }
\newcommand{\Subsection}{\subsection}
\providecommand{\nobreakdash}{\nobreak\hbox{-}}
\setlength{\emergencystretch}{2em}
\multlinegap0pt
\usepackage{amssymb}
\usepackage{mathtools}
\usepackage{tikz}
\newtheorem{proposition}{Proposition}
\newcommand{\Norm}{\operatorname{N}}
\newcommand{\Oct}{\mathbb O}
\newcommand{\QQ}{\mathbb Q}
\newcommand{\Rea}{\operatorname{Re}}
\newcommand{\Zphi}{\mathbb Z[\varphi]}
\title{Non-crystallographic systems of integers over composition algebras}
\author{Daniele Corradetti}
\date{Source: arXiv:2605.15075v1 (CC BY 4.0)}
\begin{document}
\maketitle

\noindent\emph{Source and licence.} Non-crystallographic systems of integers over composition algebras. Version arXiv:2605.15075v1 (CC BY 4.0), distributed by arXiv under the Creative Commons Attribution 4.0 International licence, http://creativecommons.org/licenses/by/4.0/, as recorded in the arXiv licence metadata for this exact version and shown on the abstract page https://arxiv.org/abs/2605.15075v1. Full licence notice: \texttt{licenses/arxiv-2605.15075-CC-BY-4.0.txt}.\par\smallskip

\noindent\textit{Editorial scope.} Counted unit: the article's proposition prop:dual-discriminant (source0.tex lines 1215--1232). The article's complete original proof of it is delivered with it (source0.tex lines 1234--1270, one \texttt{proof} environment of 37 lines). No mathematics is written, repaired or re-derived: the statement and the proof are the article's own lines, in the article's own notation, and no retained line is substituted or re-flowed.  Retained uncounted as the source-local context the proof needs: lines 997--1000.  Nothing was omitted from any retained span, so the package carries no omission record.  No retained line contains a citation, so the wrapper carries no bibliography and no bibliography entry.  No retained line contains a figure environment, an image-inclusion command or drawing code, so no image is delivered and the package declares no asset dependency.  Editorial formatting only, in addition: an AMS \texttt{amsart} preamble that restates the article's own declarations and macros used by the retained lines, namely the environments \texttt{proposition} on separate counters, together with the standard packages those lines need.  The retained environments print in this document as Proposition 1 in order of appearance, whereas the article prints the same items with the numbering of its own full document.  The complete native source of the article as distributed in the arXiv e-print for this exact version is bundled unchanged as \texttt{sources/200629/source0.tex}.
\begin{equation}
e_1=1,\qquad e_2=i,\qquad e_3=\tfrac12(1+i+j+k),\qquad e_4=\tfrac12(-1+(\varphi-1)i-\varphi j).
\label{eq:icosian-basis}
\end{equation}

\begin{proposition}[Dual-discriminant obstruction]
\label{prop:dual-discriminant}
Let
\begin{equation}
B(x,y)=\Norm(x+y)-\Norm(x)-\Norm(y)
\label{eq:norm-polar-pairing}
\end{equation}
be the polar bilinear form of the octonionic norm, and let
\begin{equation}
G_0^\#=\{x\in K\otimes_{\Zphi}G_0:\ B(x,G_0)\subset\Zphi\}.
\label{eq:g0-dual}
\end{equation}
For the icosian double \(G_0=\mathbb I\oplus\mathbb I\ell\), one has
\begin{equation}
G_0^\#=G_0.
\end{equation}
Consequently \(G_0\) has no strict norm-integral overorder obtained by adjoining a fractional ideal layer.
\end{proposition}

\begin{proof}
In the basis \((e_1,\ldots,e_4,e_1\ell,\ldots,e_4\ell)\) of \(G_0\) inherited from \eqref{eq:icosian-basis}, the Gram matrix of \(B\) is block diagonal with two identical \(4\times 4\) blocks \(G_{\mathbb I}\). A direct computation of \(B(e_i,e_j)=2\Rea(e_i\overline{e_j})\) gives
\begin{equation}
G_{\mathbb I}=
\begin{pmatrix}
2 & 0 & 1 & -1 \\
0 & 2 & 1 & \varphi-1 \\
1 & 1 & 2 & -1 \\
-1 & \varphi-1 & -1 & 2
\end{pmatrix},
\qquad
\det G_{\mathbb I}=\varphi^{2}=\varphi+1.
\label{eq:gram-icosian}
\end{equation}
Block-diagonality of the full pairing and \(\det G_{\mathbb I}=\varphi^{2}\) imply
\begin{equation}
G_{G_0}=G_{\mathbb I}\oplus G_{\mathbb I},
\qquad
\det G_{G_0}=(\det G_{\mathbb I})^{2}=\varphi^{4}=2+3\varphi.
\label{eq:gram-g0}
\end{equation}
The field norm of this determinant is
\begin{equation}
\Norm_{K/\QQ}(2+3\varphi)=(2+3\varphi)(5-3\varphi)=10+9\varphi-9\varphi^2=10+9\varphi-9(\varphi+1)=1,
\end{equation}
so \(\det G_{G_0}\) is a unit of \(\Zphi\). The adjugate \(\mathrm{adj}(G_{G_0})\) has entries in \(\Zphi\) because the cofactor of any minor of a matrix with entries in a commutative ring lies in that ring (see, e.g., Lang, \emph{Algebra}, Ch.~XIII~\S4 or the determinantal Cayley--Hamilton statement). Since \((\det G_{G_0})^{-1}\) lies in \(\Zphi\) by the unit computation above, the identity
\begin{equation}
G_{G_0}^{-1}=(\det G_{G_0})^{-1}\,\mathrm{adj}(G_{G_0})
\end{equation}
gives \(G_{G_0}^{-1}\in M_{8}(\Zphi)\). The pairing therefore identifies \(G_0\) with \(\operatorname{Hom}_{\Zphi}(G_0,\Zphi)\), and \(G_0^\#/G_0\) is trivial.

Now let \(G\supset G_0\) be an overorder of \(\Oct(K)\) with integral norm on every element. For \(x\in G\) and \(g\in G_0\),
\begin{equation}
B(x,g)=\Norm(x+g)-\Norm(x)-\Norm(g)\in\Zphi.
\end{equation}
Thus \(x\in G_0^\#=G_0\), and \(G=G_0\). \(\square\)
\end{proof}

\end{document}
